\clearpage
\section{Results}

This section evaluates the implemented hierarchical exploration
framework. Rather than describing the internal structure of Talos,
which was presented in the previous chapter, the experiments focus
on the design alternatives obtained through the exploration and on
the effect of introducing characterized hardware IP information
after the architectural exploration stage.


\subsection{Experimental Setup}

\subsubsection{Workloads}

Several workloads with different numerical representations were used
during the development and evaluation of Talos. All of them are
provided as ONNX models, since this is the workload representation
used by the ZigZag integration.

The workloads considered for the experiments are summarized in
Table~\ref{tab:workloads}.

\begin{table}[h]
\centering
\begin{tabular}{lll}
\hline
Workload & Numerical format & Use \\
\hline
AlexNet & INT16 & Complete CNN workload \\
SqueezeNet 1.0 & INT8 & Quantized integer workload \\
SqueezeNet 1.0 & FP16 & Floating-point workload \\
VGG16 & FP16 & Comparison with SAURIA \\
ResNet18 first layer & FP32 & Controlled convolution workload \\
\hline
\end{tabular}
\caption{Workloads considered during the Talos evaluation.}
\label{tab:workloads}
\end{table}

The INT16 AlexNet model follows the AlexNet topology used by Eyeriss.
It contains the five convolutional layers and the three fully
connected layers. Since the original quantized parameters are not
available, this model is used to preserve the layer dimensions and
numerical representation required by the exploration rather than as
an executable inference model.

Two versions of SqueezeNet 1.0 are also available. The INT8 version
uses quantized convolution operations with unsigned 8-bit activations,
signed 8-bit weights, and unsigned 8-bit outputs. The FP16 model
provides the same type of CNN workload using half-precision
floating-point values.

An FP16 version of VGG16 is also used for the comparison with SAURIA,
which reports results for the same network and numerical format
\cite{sauria2023}.

Finally, the first convolutional layer of ResNet18 is provided as a
smaller controlled workload. This allows complete Level 1 and Level 2
explorations to be executed with a lower simulation time while still
keeping the structure of a real convolutional layer.

Using workloads with different precisions is relevant for the Level 2
experiments because the physical PE pool also contains implementations
with different numerical formats. Talos obtains the operand
representation directly from the ONNX model and reports when the
selected PE representation does not match the workload. Talos does
not automatically quantize or convert the input network.


\subsubsection{Hardware Characterization Setup}

The physical IP characterization used for the final experiments was
performed for a TSMC 65~nm technology.

The characterized IPs were obtained using Cadence Genus after mapping
the RTL to the target standard-cell library. The common
characterization point used by the current pools is:

\begin{itemize}
    \item target reference frequency of 200~MHz;
    \item supply voltage of 1.2~V;
    \item temperature of 25~$^\circ$C;
    \item NCCOM process corner.
\end{itemize}

The power characterization provides an idle and an active operating
point for each IP. These values were obtained using the vectorless
power flow described in the previous chapter.

Memory blocks and processing elements use the same characterization
conditions so that they can later be combined in the same Level 2
implementation. The resulting areas, delays, maximum frequencies, and
power values are stored in the Talos IP pools.

Three characterized pools are considered in the experiments.

The first one,
\texttt{ip\_pool\_characterized\_65nm.yaml}, contains the integer
processing elements and the characterized memory implementations. It
is used to study architectures implemented using integer compute
units.

The second one,
\texttt{ip\_pool\_fp16\_65nm.yaml}, contains the characterized FP16
processing elements together with the same register-file and
global-buffer alternatives.

The third one,
\texttt{ip\_pool\_fp32\_65nm.yaml}, contains the characterized FP32
processing elements and the characterized memory pool.

The FP16 pool contains two processing-element alternatives, while the
FP32 pool contains three. Both use twelve register-file
implementations and three global-buffer implementations. The integer
pool contains the characterized integer processing elements and the
same memory characterization.

The different pools are intentionally kept separate. This allows the
physical exploration to be repeated using a coherent set of
processing elements for a given numerical representation, without
changing the Level 1 architectural representation.

Talos also contains an example pool and a synthetic 65~nm pool. These
were used during the development and validation of the framework, but
they are not intended to represent the final physical
characterization used to draw conclusions from the experiments.

External DRAM is treated separately from the accelerator IPs. When a
characterized DRAM is not provided, Talos uses the fixed platform
proxy described in the previous chapter. Therefore, the DRAM is not
part of the Level 2 IP selection.


\subsubsection{Exploration Configuration}

The experiments use the complete hierarchical Talos flow.

Level 1 explores the abstract architecture using NSGA-II and ZigZag.
In the current implementation, the architectural screening is
performed using energy, latency, and area at the same time. These
three objectives are used to preserve different architectural
trade-offs before physical IP selection.

The Level 1 design space contains the PE-array dimensions, local
storage capacity and bandwidth, global-buffer capacity and bandwidth,
and the PE-array dimensions served by the global buffer. The exact
definition of this genome was presented in the previous chapter.

After Level 1, a limited number of architectures from the resulting
set are transferred to Level 2. The ZigZag workload activity profile
associated with each architecture is transferred together with it, so
the mapping does not need to be regenerated during the physical
exploration.

Level 2 is evaluated using both available exploration strategies when
required. Exhaustive exploration is used when the number of compatible
IP combinations is small enough to enumerate all of them. NSGA-II can
instead be used for larger physical design spaces.

The Level 2 experiments consider area, energy, and workload
performance, both individually and in combination. When exhaustive exploration uses more than one objective, a
representative solution is selected using the logarithmic augmented
Tchebycheff score defined in the following section.

The population size, number of generations, number of architectures
transferred to Level 2, and user constraints will be reported together
with each final experiment. This avoids tying the experimental setup
to intermediate runs performed during the development of the tool.


\subsection{Hardware IP Characterization Results}

%\subsubsection{Processing Element Characterization}

% Final table generated from the characterization used in the
% experiments.
%
% Suggested columns:
% PE
% numerical format
% area
% max frequency
% delay
% active power
% idle power
% MACs/cycle
%
% Separate integer, FP16 and FP32 implementations if the resulting
% table becomes too large.
%
% Compare only the final characterization values used by the
% experiments.


%\subsubsection{Memory Characterization}

% Final table generated from the characterized memory pool.
%
% Suggested RF columns:
% memory
% ports
% capacity
% bandwidth
% area
% max frequency
% idle power
% active power
% read energy
% write energy
%
% Suggested GB columns:
% memory
% capacity
% bandwidth
% area
% max frequency
% idle power
% active power
% read/write energy
%
% Discuss scaling with capacity and bandwidth once the final
% characterization has been fixed.


\subsubsection{Available IP Design Space}

The physical design space depends on the IP pool selected for each
experiment.

The characterized integer pool contains seven PE implementations.
Six of them provide one MAC per cycle, while the OpenEye composite PE
provides two MACs per cycle and already includes the three local
register-file roles used by Talos. Most of the integer PEs use signed
INT8 arithmetic, while the TMS4517 implementation uses UINT8.

The integer pool also contains twelve register-file implementations
and three global-buffer alternatives.

The FP16 pool contains two FP16 PEs, twelve register files, and three
global buffers. The FP32 pool contains three FP32 PEs and the same
memory alternatives.

Therefore, the three characterized pools provide different compute
alternatives while reusing the same characterized memory design
space. This makes it possible to study the effect of changing the PE
implementation and numerical format without changing the available
memory IPs.

Not every IP is necessarily compatible with every abstract component.
Capacity and bandwidth requirements are taken from the Level 1
architecture and are used to reduce the compatible memory
alternatives. In addition, the workload-aware evaluation can reject a
PE or memory that cannot sustain the activity required by the mapping.

Therefore, the actual Level 2 design-space size is architecture
dependent and will be reported together with the final exploration
results.


%\subsection{Level 1 Architecture Exploration}

%\subsubsection{Explored Architectures}

% Final Level 1 results.
%
% For each workload:
% number of evaluated architectures
% number of unique architectures
% number of feasible architectures
% Pareto/front information
%
% Suggested plots:
% energy vs latency
% energy vs area
% latency vs area
%
% Use the new final runs only.


%\subsubsection{Selected Architectural Candidates}

% Architectures finally transferred to Level 2.
%
% Report:
% PE array
% RF requirements
% GB requirements
% L1 latency
% L1 energy
% L1 area
%
% Explain why the points are relevant, but do not repeat the genome
% definition.


%\subsection{IP-Based Evaluation Results}

%\subsubsection{Feasible IP Implementations}

% For every selected L1 architecture:
% theoretical compatible combinations
% evaluated combinations
% valid combinations
% rejected combinations
%
% Separate:
% compatibility rejection
% workload throughput/bandwidth rejection
% user-constraint rejection.


%\subsubsection{Impact of Processing Element Selection}

% Keep architecture and memories fixed.
%
% Compare different PEs from:
% - integer pool
% - FP16 pool
% - FP32 pool
%
% Metrics:
% area
% power
% energy/inference
% fmax
% workload feasibility.


%\subsubsection{Impact of Memory Selection}

% Keep architecture and PE fixed.
%
% Change RF / GB implementation.
%
% Compare:
% area
% power
% energy
% bandwidth
% accesses/cycle
% feasibility.

%\subsubsection{Workload-Aware PPA Results}

% Main Level 2 result table.
%
% Suggested columns:
% workload
% L1 architecture
% pool
% PE
% RFs
% GB
% area
% energy/inference
% average power
% workload latency
% workload throughput
% physical fmax
% balanced score if applicable.


\subsection{Target-Oriented Architecture Selection}

A complete target sweep was performed for the INT16 AlexNet workload
in order to study how the preferred architecture changes with the
optimization objective. The sweep considered single-objective targets
for energy, area, and performance, as well as the main multi-objective
combinations.

For the single-objective cases, the winning architecture is the one
with the minimum value of the corresponding objective. For the
multi-objective cases, the objectives are expressed as logarithmic
relative deviations from the ideal point,

\[
r_i(x)=
\ln\left(
\frac{f_i(x)}
{f_i^{\min}}
\right),
\]

and candidates are ranked using an augmented Tchebycheff score,

\[
S(x)=
100\left[
\max_i r_i(x)
+
0.05\sum_i r_i(x)
\right].
\]

This formulation emphasizes the worst relative degradation with
respect to the ideal value while using the sum of deviations as a
secondary discrimination term. The ideal values are computed over the
feasible solutions considered in the corresponding sweep. When
architectures, seeds, or runs are compared, Talos recalculates one
common ideal over the complete feasible set being compared. Local
scores calculated inside an individual architecture are not used for
cross-architecture selection.

The Level 2 energy objective is \texttt{workload\_energy\_j}, which
includes external DRAM energy. On-chip energy is computed as workload
energy minus DRAM energy and is shown only for reporting and
comparison.

Table~\ref{tab:alexnet_target_winners} summarizes the winning Level~1
architecture for each target. The energy, area, and latency columns
correspond to the selected feasible Level~2 implementation associated
with the architecture. For the performance-only target, several
Level~2 implementations can sustain the same Level~1 mapping and
therefore tie in workload latency; in that case, secondary metrics are
used only to choose a representative physical implementation.

\begin{table*}[t]
\centering
\caption{Winning architectures for the INT16 AlexNet target sweep
using the augmented Tchebycheff score for multi-objective selection.
Energy values exclude external DRAM energy.}
\label{tab:alexnet_target_winners}

\setlength{\tabcolsep}{4pt}
\renewcommand{\arraystretch}{1.15}

\resizebox{0.98\textwidth}{!}{%
\begin{tabular}{@{}lcccccccc@{}}
\hline
\textbf{Target} &
\textbf{Arch.} &
\textbf{PE array} &
\makecell{\textbf{RF size}\\\textbf{(bit)}} &
\makecell{\textbf{RF BW}\\\textbf{(bit/cycle)}} &
\makecell{\textbf{GB size}\\\textbf{(Kib)}} &
\makecell{\textbf{GB BW}\\\textbf{(bit/cycle)}} &
\makecell{\textbf{On-chip Energy}\\\textbf{(mJ/inf.)$^{a}$}} &
\makecell{\textbf{Area / Latency}\\\textbf{(mm$^2$ / ms)}} \\
\hline

Energy &
144 &
$32\times 8$ &
128 &
16 &
128 &
1024 &
44.99 &
3.251 / 58.95 \\

Energy + Area &
46 &
$4\times 8$ &
64 &
8 &
32 &
128 &
88.29 &
0.512 / 810.52 \\

Area &
146 &
$4\times 4$ &
256 &
16 &
8 &
64 &
65.39 &
0.2653 / 1129.8 \\

Performance &
67 &
$32\times 32$ &
256 &
256 &
64 &
1024 &
4627.69 &
1491.556 / 16.58 \\

Area + Performance &
108 &
$16\times 8$ &
256 &
32 &
32 &
256 &
50.14 &
1.780 / 116.59 \\

Energy + Performance &
104 &
$32\times 32$ &
64 &
16 &
128 &
1024 &
47.81 &
8.832 / 21.08 \\

Energy + Area + Performance &
108 &
$16\times 8$ &
256 &
32 &
32 &
256 &
50.14 &
1.780 / 116.59 \\

\hline
\end{tabular}%
}

\vspace{1mm}
\begin{minipage}{0.96\textwidth}
\footnotesize
$^{a}$ The reported energy corresponds to the on-chip contribution,
computed as
$E_{\mathrm{on-chip}} = E_{\mathrm{workload}} - E_{\mathrm{DRAM}}$.
The architectures were selected using the original sweep energy
objective, which considers total workload energy including external
DRAM. DRAM energy is removed here only for reporting and comparison
purposes.
\end{minipage}

\end{table*}


Table with optimizing at level 1.

\begin{table*}[t]
\centering
\caption{Winning architectures for the legacy INT16 AlexNet target
sweep using an independent ZigZag-based Level~1 search for each
optimization target and the augmented Tchebycheff score for
multi-objective selection.}
\label{tab:alexnet_legacy_target_winners}

\setlength{\tabcolsep}{4pt}
\renewcommand{\arraystretch}{1.15}

\resizebox{0.98\textwidth}{!}{%
\begin{tabular}{@{}lcccccccc@{}}
\hline
\textbf{Target} &
\textbf{Arch.} &
\textbf{PE array} &
\makecell{\textbf{RF size}\\\textbf{(bit)}} &
\makecell{\textbf{RF BW}\\\textbf{(bit/cycle)}} &
\makecell{\textbf{GB size}\\\textbf{(Kib)}} &
\makecell{\textbf{GB BW}\\\textbf{(bit/cycle)}} &
\makecell{\textbf{On-chip Energy}\\\textbf{(mJ/inf.)$^{a}$}} &
\makecell{\textbf{Area / Latency}\\\textbf{(mm$^2$ / ms)}} \\
\hline

Energy &
5 &
$8\times 32$ &
128 &
16 &
128 &
1024 &
49.69 &
3.251 / 66.19 \\

Energy + Area &
3 &
$4\times 8$ &
256 &
8 &
8 &
64 &
103.53 &
0.445 / 1048.94 \\

Area &
0 &
$4\times 4$ &
64 &
8 &
8 &
64 &
68.80 &
0.171 / 1868.89 \\

Performance &
0 &
$32\times 32$ &
512 &
256 &
128 &
64 &
4597.68 &
1491.556 / 16.56 \\

Area + Performance &
28 &
$16\times 8$ &
256 &
32 &
32 &
256 &
48.74 &
1.780 / 113.07 \\

Energy + Performance &
24 &
$16\times 32$ &
512 &
16 &
8 &
256 &
120.54 &
20.733 / 27.35 \\

Energy + Area + Performance &
50 &
$16\times 4$ &
64 &
8 &
32 &
128 &
80.93 &
0.682 / 522.69 \\

\hline
\end{tabular}%
}

\vspace{1mm}
\begin{minipage}{0.96\textwidth}
\footnotesize

$^{a}$ The reported energy corresponds to the on-chip contribution,
computed as
$E_{\mathrm{on-chip}} =
E_{\mathrm{workload}} - E_{\mathrm{DRAM}}$.
The architectures were selected using total workload energy including
external DRAM. DRAM energy is removed here only for reporting and
comparison purposes.

Unlike the common-handoff sweep, Level~1 is executed independently
for each optimization target in this experiment. Consequently,
architecture indices are local to each target and the same architecture
index in two different rows does not necessarily represent the same
hardware configuration.

\end{minipage}
\end{table*}

\begin{table*}[t]
\centering
\caption{Best Talos implementations obtained for the INT16 AlexNet
workload used for comparison with Eyeriss V2}
\label{tab:eyeriss_int16_v2_results}

\setlength{\tabcolsep}{3.2pt}
\renewcommand{\arraystretch}{1.15}

\resizebox{\textwidth}{!}{%
\begin{tabular}{@{}ccccccclccccccc@{}}
\hline
\textbf{Obj.} &
\makecell{\textbf{Best}\\\textbf{run}} &
\makecell{\textbf{Energy}\\\textbf{[J]}} &
\makecell{\textbf{Latency}\\\textbf{[s]}} &
\makecell{\textbf{Inf.}\\\textbf{[s$^{-1}$]}} &
\makecell{\textbf{Area}\\\textbf{[mm$^2$]}} &
\makecell{\textbf{Power}\\\textbf{[W]}} &
\textbf{PE} &
\makecell{\textbf{PE}\\\textbf{array}} &
\makecell{\textbf{RF size}\\\textbf{[bit]}} &
\makecell{\textbf{RF BW}\\\textbf{[bit/cycle]}} &
\makecell{\textbf{GB size}\\\textbf{[bit]}} &
\makecell{\textbf{GB BW}\\\textbf{[bit/cycle]}} &
\makecell{\textbf{L1}\\\textbf{arch.}} &
\makecell{\textbf{L2}\\\textbf{sol.}} \\
\hline

E &
seed\_24 &
0.112869 &
0.0316643 &
31.5813 &
5.111114 &
3.564564 &
\texttt{pe\_tiny\_xpu\_int16} &
$16\times32$ &
64 &
16 &
131072 &
1024 &
39 &
2555 \\

E+A &
seed\_22 &
0.244724 &
0.810516 &
1.23378 &
0.512155 &
0.301937 &
\texttt{pe\_tiny\_xpu\_int16} &
$8\times4$ &
64 &
8 &
32768 &
128 &
9 &
8796 \\

A &
large\_pool &
0.450138 &
1.240105 &
0.806383 &
0.201889 &
0.362984 &
\texttt{pe\_tiny\_xpu\_int16} &
$4\times4$ &
128 &
16 &
8192 &
64 &
42 &
7665 \\

P &
seed\_25 &
4.708420 &
0.0165771 &
60.3241 &
1493.550121 &
284.031252 &
\texttt{pe\_sauria\_int16} &
$32\times32$ &
256 &
256 &
65536 &
1024 &
8 &
0 \\

A+P &
seed\_21 &
0.198234 &
0.111362 &
8.97969 &
1.700708 &
1.780077 &
\texttt{pe\_tiny\_xpu\_int16} &
$16\times16$ &
64 &
8 &
16384 &
256 &
69 &
8796 \\

E+P &
seed\_24 &
0.124780 &
0.0197103 &
50.7350 &
8.831982 &
6.330738 &
\texttt{pe\_tiny\_xpu\_int16} &
$32\times32$ &
64 &
16 &
65536 &
1024 &
51 &
2555 \\

E+A+P &
seed\_21 &
0.198234 &
0.111362 &
8.97969 &
1.700708 &
1.780077 &
\texttt{pe\_tiny\_xpu\_int16} &
$16\times16$ &
64 &
8 &
16384 &
256 &
69 &
8796 \\

\hline
\end{tabular}%
}

\vspace{1mm}
\begin{minipage}{0.97\textwidth}
\footnotesize
The workload corresponds to the INT16 AlexNet configuration used for
comparison with Eyeriss. E, A, and P denote energy per inference,
area, and workload performance, respectively. The reported
architecture and solution indices identify the selected Level~1
architecture and Level~2 physical implementation.
\end{minipage}
\end{table*}

\begin{table*}[t]
\centering
\caption{Representative architectures selected from the exhaustive
INT16 Eyeriss AlexNet convolution design space using an augmented
Tchebycheff score for multi-objective selection. (4batch)}
\label{tab:eyeriss_int16_target_winners}

\setlength{\tabcolsep}{4pt}
\renewcommand{\arraystretch}{1.15}

\resizebox{0.98\textwidth}{!}{%
\begin{tabular}{@{}lcccccccc@{}}
\hline
\textbf{Target} &
\textbf{Arch.} &
\textbf{PE array} &
\makecell{\textbf{RF size}\\\textbf{(bit)}} &
\makecell{\textbf{RF BW}\\\textbf{(bit/cycle)}} &
\makecell{\textbf{GB size}\\\textbf{(Kib)}} &
\makecell{\textbf{GB BW}\\\textbf{(bit/cycle)}} &
\makecell{\textbf{On-chip Energy}\\\textbf{(mJ/image)$^{a}$}} &
\makecell{\textbf{Area / Latency}\\\textbf{(mm$^2$ / ms)}} \\
\hline

Energy &
14 &
$32\times 8$ &
128 &
16 &
128 &
512 &
44.06 &
3.251 / 242.15 \\

Energy + Area &
41 &
$8\times 4$ &
128 &
16 &
8 &
64 &
45.59 &
0.575 / 1472.64 \\

Area &
27 &
$4\times 4$ &
256 &
32 &
8 &
64 &
45.54 &
0.2653 / 3118.17 \\

Performance &
22 &
$32\times 32$ &
512 &
128 &
32 &
128 &
550.90 &
386.255 / 29.69 \\

Area + Performance &
14 &
$32\times 8$ &
128 &
16 &
128 &
512 &
44.06 &
3.251 / 242.15 \\

Energy + Performance &
33 &
$32\times 32$ &
256 &
128 &
16 &
256 &
78.84 &
46.633 / 30.20 \\

Energy + Area + Performance &
14 &
$32\times 8$ &
128 &
16 &
128 &
512 &
44.06 &
3.251 / 242.15 \\

\hline
\end{tabular}%
}

\vspace{1mm}
\begin{minipage}{0.96\textwidth}
\footnotesize
$^{a}$ The reported energy corresponds to the on-chip contribution,
computed as
$E_{\mathrm{on-chip}} =
E_{\mathrm{workload}} - E_{\mathrm{DRAM}}$.
The Eyeriss-compatible AlexNet workload uses a batch size of four;
therefore, the reported value is
$E_{\mathrm{on-chip}}/4$ and is expressed in mJ/image.
The architectures are selected using total workload energy including
external DRAM. DRAM energy is removed only from the reported values.
Latency corresponds to the execution of the complete batch.
\end{minipage}

\end{table*}
% NOTE: The FP16 and FP32 tables below still use the winners from
% their original result post-processing. Recomputing them with the
% augmented Tchebycheff score requires the corresponding candidate CSVs.
% They are intentionally left unchanged here rather than inferring new
% winners from the already-selected rows.
\begin{table*}[t]
\centering
\caption{Winning architectures for the FP16 target sweep using an
independent ZigZag-based Level~1 search for each optimization target.}
\label{tab:fp16_target_winners}

\setlength{\tabcolsep}{4pt}
\renewcommand{\arraystretch}{1.15}

\resizebox{0.98\textwidth}{!}{%
\begin{tabular}{@{}lcccccccc@{}}
\hline
\textbf{Target} &
\textbf{Arch.} &
\textbf{PE array} &
\makecell{\textbf{RF size}\\\textbf{(bit)}} &
\makecell{\textbf{RF BW}\\\textbf{(bit/cycle)}} &
\makecell{\textbf{GB size}\\\textbf{(Kib)}} &
\makecell{\textbf{GB BW}\\\textbf{(bit/cycle)}} &
\makecell{\textbf{On-chip Energy}\\\textbf{(mJ/inf.)$^{a}$}} &
\makecell{\textbf{Area / Latency}\\\textbf{(mm$^2$ / ms)}} \\
\hline

Energy &
8 &
$32\times32$ &
512 &
16 &
32 &
256 &
37.02 &
32.070 / 4.97 \\

Energy + Area &
13 &
$4\times4$ &
128 &
32 &
32 &
256 &
38.45 &
0.546 / 304.39 \\

Area &
13 &
$4\times4$ &
128 &
32 &
32 &
256 &
38.45 &
0.546 / 304.39 \\

Performance &
0 &
$32\times32$ &
64 &
64 &
64 &
512 &
1091.94 &
1444.727 / 3.93 \\

Area + Performance &
30 &
$16\times4$ &
128 &
16 &
64 &
1024 &
37.70 &
1.953 / 86.83 \\

Energy + Performance &
8 &
$32\times32$ &
512 &
16 &
32 &
256 &
37.02 &
32.070 / 4.97 \\

Energy + Area + Performance &
30 &
$16\times4$ &
128 &
16 &
64 &
1024 &
37.70 &
1.953 / 86.83 \\

\hline
\end{tabular}%
}

\vspace{1mm}
\begin{minipage}{0.96\textwidth}
\footnotesize

$^{a}$ The reported energy corresponds to the on-chip contribution,
computed as
$E_{\mathrm{on-chip}} =
E_{\mathrm{workload}} - E_{\mathrm{DRAM}}$.
The architectures are selected using the original optimization
objectives, where energy corresponds to total workload energy
including external DRAM. DRAM energy is removed here only for
reporting and comparison purposes.

Level~1 is executed independently for each optimization target.
Consequently, architecture indices are local to each target and the
same index in different rows does not necessarily identify the same
architecture. For the performance-only target, multiple Level~2
implementations achieve the same minimum workload latency; the
reported physical metrics correspond to the lowest-energy
implementation among the tied solutions.

\end{minipage}

\end{table*}

\begin{table*}[t]
\centering
\caption{Winning architectures for the FP32 target sweep using an
independent ZigZag-based Level~1 search for each optimization target.}
\label{tab:fp32_target_winners}

\setlength{\tabcolsep}{4pt}
\renewcommand{\arraystretch}{1.15}

\resizebox{0.98\textwidth}{!}{%
\begin{tabular}{@{}lcccccccc@{}}
\hline
\textbf{Target} &
\textbf{Arch.} &
\textbf{PE array} &
\makecell{\textbf{RF size}\\\textbf{(bit)}} &
\makecell{\textbf{RF BW}\\\textbf{(bit/cycle)}} &
\makecell{\textbf{GB size}\\\textbf{(Kib)}} &
\makecell{\textbf{GB BW}\\\textbf{(bit/cycle)}} &
\makecell{\textbf{On-chip Energy}\\\textbf{(mJ/inf.)$^{a}$}} &
\makecell{\textbf{Area / Latency}\\\textbf{(mm$^2$ / ms)}} \\
\hline

Energy &
53 &
$16\times16$ &
128 &
128 &
128 &
1024 &
15.43 &
12.649 / 4.63 \\

Energy + Area &
45 &
$8\times4$ &
1024 &
32 &
32 &
256 &
17.73 &
1.750 / 44.16 \\

Area &
27 &
$4\times4$ &
512 &
64 &
16 &
256 &
25.14 &
0.794 / 128.96 \\

Performance &
6 &
$32\times32$ &
128 &
128 &
64 &
512 &
366.06 &
1468.647 / 1.32 \\

Area + Performance &
53 &
$16\times16$ &
128 &
128 &
128 &
1024 &
15.43 &
12.649 / 4.63 \\

Energy + Performance &
30 &
$32\times32$ &
1024 &
32 &
32 &
1024 &
25.12 &
89.523 / 1.32 \\

Energy + Area + Performance &
53 &
$16\times16$ &
128 &
128 &
128 &
1024 &
15.43 &
12.649 / 4.63 \\

\hline
\end{tabular}%
}

\vspace{1mm}
\begin{minipage}{0.96\textwidth}
\footnotesize

$^{a}$ The reported energy corresponds to the on-chip contribution,
computed as
$E_{\mathrm{on-chip}} =
E_{\mathrm{workload}} - E_{\mathrm{DRAM}}$.
The architectures are selected using the original optimization
objectives, where energy corresponds to total workload energy
including external DRAM. DRAM energy is removed here only for
reporting and comparison purposes.

Level~1 is executed independently for each optimization target.
Consequently, architecture indices are local to each target and the
same index in different rows does not necessarily identify the same
architecture. For the performance-only target, multiple Level~2
implementations achieve the same minimum workload latency; the
reported physical metrics correspond to the lowest-energy
implementation among the tied solutions.

\end{minipage}

\end{table*}



The results show that the preferred architecture changes substantially
when optimizing the individual extremes. Minimum area favours a small
$4\times4$ PE array, while minimum energy favours a $32\times8$
array and maximum performance favours a $32\times32$ array. This
confirms that the architectural parameters cannot be selected
independently of the optimization target.

The multi-objective selections expose different compromises.
Architecture~108 is selected for both area--performance and the
combined energy--area--performance objective. It uses a
$16\times8$ PE array and reaches approximately 116.59~ms per
inference with 1.78~mm$^2$ of area and 50.14~mJ of reported on-chip
energy. In contrast, architecture~104 is selected for the
energy--performance objective. Its larger $32\times32$ array reduces
latency to approximately 21.08~ms, at the cost of increasing area to
8.83~mm$^2$, while retaining a similar on-chip energy of 47.81~mJ.
This illustrates how introducing area into the objective favours a
smaller implementation even when this requires a substantial latency
trade-off.

For the area-only target, architectures~31 and~146 obtain the same
minimum physical area of approximately 0.2653~mm$^2$. Table~
\ref{tab:alexnet_target_winners} reports architecture~146 as the
representative solution. The area-only result should therefore be
interpreted as a tie at the primary objective rather than as a unique
global winner.

For the performance-only target, architecture~67 reaches approximately
16.58~ms per inference, or 60.3~inferences/s. Several Level~2
implementations can sustain the same Level~1 mapping and therefore tie
in workload latency. The reported physical metrics correspond to one
representative implementation of the winning architectural
configuration.

\subsection{Comparison with Published Accelerators}

To give some context to the implementations generated by Talos, one
representative accelerator configuration is compared with DNN
accelerators reported in the literature.

This comparison is intended as a reference rather than as a direct
benchmark. The accelerators use different technologies, numerical
precisions, workloads, operating points, and architectures.

\begin{table*}[t]
\centering
\caption{Comparison of representative DNN accelerators and the
Talos design selected for the combined energy, area, and performance
objective.}
\label{tab:dnn_accelerator_comparison}

\setlength{\tabcolsep}{4pt}
\renewcommand{\arraystretch}{1.15}

\resizebox{0.98\textwidth}{!}{%
\begin{tabular}{@{}lcccccc@{}}
\hline
\textbf{Accelerator} &
\textbf{Technology} &
\textbf{Precision} &
\makecell{\textbf{Energy /}\\\textbf{inf.}} &
\textbf{Area} &
\textbf{Power} &
\textbf{Performance} \\
\hline

TinyVers \cite{tinyvers} &
22\,nm FDX &
INT2/4/8 &
\makecell{0.56--17.3\\$\mu$J$^{a}$} &
6.25\,mm$^2$ &
\makecell{1.7\,$\mu$W--\\20\,mW} &
\makecell{up to\\17.6\,GOPS} \\

Eyeriss \cite{chen2017eyeriss} &
\makecell{65\,nm\\LP CMOS} &
INT16 &
8.01\,mJ$^{b}$ &
12.25\,mm$^2$ &
278\,mW &
34.7 inf/s \\

ENVISION \cite{envision,chen2019eyerissv2} &
\makecell{28\,nm\\FDSOI} &
4/8/16 bit &
0.936\,mJ$^{b}$ &
1.87\,mm$^2$ &
7.6--290\,mW &
$\sim$47 inf/s \\

UNPU \cite{lee2018unpu,chen2019eyerissv2} &
\makecell{65\,nm\\CMOS} &
1--16 bit &
0.911\,mJ$^{b}$ &
16\,mm$^2$ &
3.2--297\,mW &
\makecell{691.2\,GOPS (8 bit)\\
7.37\,TOPS (1 bit)} \\

Eyeriss v2 \cite{chen2019eyerissv2} &
\makecell{65\,nm\\CMOS} &
INT8 &
\makecell{1.50\,mJ$^{c}$\\
0.391\,mJ$^{d}$} &
\makecell{2695k NAND2\\
+ 246\,kB SRAM} &
\makecell{$\sim$419--\\574\,mW$^{e}$} &
\makecell{278.7 inf/s\\
1470.6 inf/s} \\

SAURIA \cite{sauria2023} &
\makecell{22\,nm\\GF22FDX} &
FP16 &
--- &
0.897\,mm$^2$ &
258\,mW &
\makecell{$\sim$6.10 inf/s$^{f}$\\
284\,GFLOP/s peak} \\

\textbf{Talos (E+A+P)} &
\makecell{\textbf{65\,nm}\\\textbf{CMOS}} &
\textbf{INT16} &
\textbf{50.14\,mJ$^{g}$} &
\textbf{1.78\,mm$^2$} &
\textbf{0.430\,W$^{g}$} &
\textbf{8.58 inf/s} \\

\hline
\end{tabular}%
}

\vspace{1mm}
\begin{minipage}{0.92\textwidth}
\footnotesize

$^{a}$ TinyVers energy depends on the workload:
0.555\,$\mu$J for OC-SVM, 2.12\,$\mu$J for TCN keyword spotting,
6.27\,$\mu$J for CAE machine monitoring, and 17.33\,$\mu$J for
ResNet-8 on CIFAR-10 \cite{tinyvers}.

$^{b}$ Energy per inference corresponds to the convolutional layers
of AlexNet. The Eyeriss value is derived from its reported power and
throughput, while the ENVISION and UNPU values are taken from the
comparison reported in \cite{chen2019eyerissv2}.

$^{c}$ Eyeriss v2 energy per inference for AlexNet.

$^{d}$ Eyeriss v2 energy per inference for MobileNet.

$^{e}$ Approximate Eyeriss v2 power derived from the reported
throughput and inference efficiency as
$P = (\mathrm{inf/s})/(\mathrm{inf/J})$
\cite{chen2019eyerissv2}.

$^{f}$ SAURIA reports a VGG-16 execution time of 164\,ms using its
on-the-fly data feeder and assuming 6.4\,GB/s of external DRAM
bandwidth, corresponding to approximately 6.10\,inferences/s.
The reported 258\,mW corresponds to accelerator power during
computation. The paper does not directly report energy per VGG-16
inference \cite{sauria2023}.

$^{g}$ Talos results correspond to the Level~2 implementation
selected for the combined energy, area, and performance objective
(architecture~108). Energy and power exclude the external DRAM
contribution. The on-chip energy is computed as
$E_{\mathrm{on-chip}} =
E_{\mathrm{workload}} - E_{\mathrm{DRAM}}$,
and the reported power corresponds to the resulting average on-chip
power over the workload execution. The Talos workload is the complete
INT16 AlexNet model used in the exploration.

Cross-accelerator comparisons should be interpreted as indicative
rather than direct benchmarks. The designs differ in CMOS technology,
numerical precision, workload definition, operating conditions, and
reported performance metric. In particular, some published AlexNet
energy values include only the convolutional layers, whereas the Talos
results correspond to the complete workload used in the exploration.
\end{minipage}
\end{table*}

\begin{table*}[t]
\centering
\caption{Best Talos implementations obtained for the FP16 VGG16
workload under the different Level 2 objective combinations.}
\label{tab:vgg16_fp16_results}

\setlength{\tabcolsep}{3.2pt}
\renewcommand{\arraystretch}{1.15}

\resizebox{\textwidth}{!}{%
\begin{tabular}{@{}ccccccclccccccc@{}}
\hline
\textbf{Obj.} &
\makecell{\textbf{Best}\\\textbf{run}} &
\makecell{\textbf{Energy}\\\textbf{[J]}} &
\makecell{\textbf{Latency}\\\textbf{[s]}} &
\makecell{\textbf{Inf.}\\\textbf{[s$^{-1}$]}} &
\makecell{\textbf{Area}\\\textbf{[mm$^2$]}} &
\makecell{\textbf{Power}\\\textbf{[W]}} &
\textbf{PE} &
\makecell{\textbf{PE}\\\textbf{array}} &
\makecell{\textbf{RF size}\\\textbf{[bit]}} &
\makecell{\textbf{RF BW}\\\textbf{[bit/cycle]}} &
\makecell{\textbf{GB size}\\\textbf{[KiB]}} &
\makecell{\textbf{GB BW}\\\textbf{[bit/cycle]}} &
\makecell{\textbf{L1}\\\textbf{arch.}} &
\makecell{\textbf{L2}\\\textbf{sol.}} \\
\hline

E &
seed\_8 &
1.20161 &
0.202291 &
4.94338 &
14.4504 &
5.94003 &
\texttt{pe\_fpnew\_fp16\_fma} &
$32\times32$ &
256 &
16 &
128 &
512 &
75 &
804 \\

E+A &
seed\_2 &
2.68957 &
11.0469 &
0.0905232 &
0.483009 &
0.243469 &
\texttt{pe\_fpnew\_fp16\_fma} &
$4\times4$ &
64 &
16 &
32 &
256 &
46 &
3110 \\

A &
seed\_10 &
5.78485 &
26.2478 &
0.0380984 &
0.226259 &
0.220394 &
\texttt{pe\_fpnew\_fp16\_fma} &
$4\times4$ &
128 &
8 &
8 &
64 &
10 &
4665 \\

P &
seed\_4 &
4.34541 &
0.161779 &
6.18126 &
114.82 &
26.8601 &
\texttt{pe\_sauria\_fp16} &
$32\times32$ &
1024 &
256 &
128 &
1024 &
17 &
0 \\

A+P &
seed\_2 &
2.12187 &
1.79265 &
0.557832 &
1.97488 &
1.18365 &
\texttt{pe\_fpnew\_fp16\_fma} &
$16\times8$ &
256 &
32 &
32 &
256 &
12 &
1608 \\

E+P &
seed\_3 &
1.32083 &
0.1996 &
5.01002 &
14.4504 &
6.6174 &
\texttt{pe\_fpnew\_fp16\_fma} &
$32\times32$ &
128 &
32 &
64 &
1024 &
56 &
804 \\

E+A+P &
large\_pool &
1.33646 &
0.805047 &
1.24216 &
3.64061 &
1.6601 &
\texttt{pe\_fpnew\_fp16\_fma} &
$32\times8$ &
128 &
16 &
128 &
1024 &
93 &
1555 \\

\hline
\end{tabular}%
}

\vspace{1mm}
\begin{minipage}{0.97\textwidth}
\footnotesize
The workload corresponds to VGG16 adapted to ONNX using FP16
operands. Level 1 is explored using energy, latency, and area
simultaneously. The first column indicates the Level 2 objective
combination, where E, A, and P denote energy per inference, area, and
workload performance, respectively.

The reported implementations were selected from ten independent runs
using the standard exploration configuration and an additional run
using a larger exploration pool. The common ideal used for the
multi-objective comparison was recalculated over the complete feasible
set from all eleven runs.
As a result, the E+A+P case selects the implementation obtained from
the larger exploration pool.

\end{minipage}
\end{table*}
\begin{table*}[t]
\centering
\caption{Best Talos implementations obtained for the INT8 MobileNet
workload under the different Level~2 objective combinations.}
\label{tab:mobilenet_int8_results}

\setlength{\tabcolsep}{3.5pt}
\renewcommand{\arraystretch}{1.15}

\resizebox{0.98\textwidth}{!}{%
\begin{tabular}{@{}lccccc@{}}
\hline
\textbf{Case} &
\makecell{\textbf{Energy}\\\textbf{[J]}} &
\makecell{\textbf{Area}\\\textbf{[mm$^2$]}} &
\makecell{\textbf{Latency}\\\textbf{[s]}} &
\makecell{\textbf{Performance}\\\textbf{[inf/s]}} &
\makecell{\textbf{Power}\\\textbf{[W]}} \\
\hline

E &
0.112869445 &
5.111114 &
0.031664310 &
31.5813 &
3.564564 \\

A &
0.446843088 &
0.265306 &
1.071744580 &
0.9331 &
0.416931 \\

E+A &
0.244724465 &
0.512155 &
0.810516235 &
1.2338 &
0.301937 \\

P &
4.708420140 &
1493.550121 &
0.016577120 &
60.3241 &
284.031252 \\

A+P &
0.198233640 &
1.700708 &
0.111362410 &
8.9797 &
1.780077 \\

E+P &
0.124780467 &
8.831982 &
0.019710255 &
50.7350 &
6.330738 \\

E+A+P &
0.198233640 &
1.700708 &
0.111362410 &
8.9797 &
1.780077 \\

\hline
\end{tabular}%
}

\vspace{1mm}
\begin{minipage}{0.96\textwidth}
\footnotesize
The workload corresponds to MobileNet using INT8 operands.
E, A, and P denote energy per inference, area, and workload
performance, respectively.
\end{minipage}
\end{table*}


\subsection{TABLAS A ENSEÑAR}

\begin{table*}[t]
\centering
\caption{Talos design configurations obtained for an INT16 AlexNet
workload under different optimization objectives.}
\label{tab:alexnet_int16_talos}

\setlength{\tabcolsep}{3.2pt}
\renewcommand{\arraystretch}{1.15}

\resizebox{0.99\textwidth}{!}{%
\begin{tabular}{@{}lcccccccc@{}}
\hline
\textbf{Configuration} &
\textbf{Technology} &
\textbf{Precision} &
\makecell{\textbf{On-chip Energy}\\\textbf{(mJ/inf.)$^{a}$}} &
\makecell{\textbf{Energy incl. DRAM}\\\textbf{(mJ/inf.)$^{b}$}} &
\textbf{Area} &
\makecell{\textbf{On-chip}\\\textbf{Power$^{a}$}} &
\makecell{\textbf{Total}\\\textbf{Power$^{b}$}} &
\makecell{\textbf{Latency}\\\textbf{(cycles)$^{c}$}} \\
\hline

Talos (E) &
65\,nm CMOS &
INT16 &
39.757 &
112.869 &
5.111\,mm$^2$ &
1.256\,W &
3.565\,W &
6332862 \\

Talos (E+A) &
65\,nm CMOS &
INT16 &
88.076 &
244.724 &
0.512\,mm$^2$ &
0.109\,W &
0.302\,W &
162103247 \\

Talos (A) &
65\,nm CMOS &
INT16 &
55.723 &
450.138 &
0.202\,mm$^2$ &
0.045\,W &
0.363\,W &
248021034 \\

Talos (P) &
65\,nm CMOS &
INT16 &
4650.562 &
4708.420 &
1493.550\,mm$^2$ &
280.541\,W &
284.031\,W &
3315424 \\

Talos (A+P) &
65\,nm CMOS &
INT16 &
51.819 &
198.911 &
1.701\,mm$^2$ &
0.465\,W &
1.786\,W &
22272482 \\

Talos (E+P) &
65\,nm CMOS &
INT16 &
45.035 &
124.780 &
8.832\,mm$^2$ &
2.285\,W &
6.331\,W &
3942051 \\

Talos (E+A+P) &
65\,nm CMOS &
INT16 &
51.142 &
198.234 &
1.701\,mm$^2$ &
0.459\,W &
1.780\,W &
22272482 \\

\hline
\end{tabular}%
}

\vspace{2mm}
\begin{minipage}{0.99\textwidth}
\footnotesize

$^{a}$ On-chip energy includes the selected processing elements and
on-chip memory hierarchy, excluding the modeled external DRAM
contribution. On-chip power is computed over the complete workload
execution.

$^{b}$ Total energy and power include the modeled external DRAM
contribution. The total workload energy is the quantity used by Talos
when energy is included as an optimization objective.

$^{c}$ Latency is expressed in workload cycles assuming the 200\,MHz
operating frequency used for the Talos characterization.

The workload corresponds to AlexNet using 16-bit integer operands.
AlexNet is a standard convolutional neural network architecture;
however, INT16 is not part of the original AlexNet specification and
should be understood here as the numerical representation selected
for this Talos experiment.

\end{minipage}

\end{table*}

\begin{table*}[t]
\centering
\caption{Comparison of the Talos FP16 VGG-16 configurations selected
under different optimization objectives and the SAURIA accelerator.}
\label{tab:vgg16_fp16_sauria_comparison}

\setlength{\tabcolsep}{3.2pt}
\renewcommand{\arraystretch}{1.15}

\resizebox{0.99\textwidth}{!}{%
\begin{tabular}{@{}lcccccccc@{}}
\hline
\textbf{Configuration} &
\textbf{Technology} &
\textbf{Precision} &
\makecell{\textbf{On-chip Energy}\\\textbf{(mJ/inf.)$^{a}$}} &
\makecell{\textbf{Energy incl. DRAM}\\\textbf{(mJ/inf.)$^{b}$}} &
\textbf{Area} &
\makecell{\textbf{On-chip}\\\textbf{Power$^{a}$}} &
\makecell{\textbf{Total}\\\textbf{Power$^{b}$}} &
\textbf{Performance} \\
\hline

Talos (E) &
65\,nm CMOS &
FP16 &
763.4 &
1201.6 &
14.450\,mm$^2$ &
3.774\,W &
5.940\,W &
4.94 inf/s \\

Talos (E+A) &
65\,nm CMOS &
FP16 &
1259.4 &
2689.6 &
0.483\,mm$^2$ &
0.114\,W &
0.243\,W &
0.091 inf/s \\

Talos (A) &
65\,nm CMOS &
FP16 &
1534.2 &
5784.9 &
0.226\,mm$^2$ &
0.058\,W &
0.220\,W &
0.038 inf/s \\

Talos (P) &
65\,nm CMOS &
FP16 &
4239.2 &
4345.4 &
114.820\,mm$^2$ &
26.203\,W &
26.860\,W &
6.18 inf/s \\

Talos (A+P) &
65\,nm CMOS &
FP16 &
927.8 &
2121.9 &
1.975\,mm$^2$ &
0.518\,W &
1.184\,W &
0.558 inf/s \\

Talos (E+P) &
65\,nm CMOS &
FP16 &
790.5 &
1320.8 &
14.450\,mm$^2$ &
3.961\,W &
6.617\,W &
5.01 inf/s \\

Talos (E+A+P) &
65\,nm CMOS &
FP16 &
762.3 &
1336.5 &
3.641\,mm$^2$ &
0.947\,W &
1.660\,W &
1.24 inf/s \\

\hline

SAURIA \cite{sauria2023} &
\makecell{22\,nm\\GF22FDX} &
FP16 &
--$^{c}$ &
--$^{c}$ &
0.897\,mm$^2$ &
0.258\,W$^{d}$ &
--$^{c}$ &
284\,GFLOP/s peak$^{d}$ \\

\hline
\end{tabular}%
}

\vspace{1mm}
\begin{minipage}{0.97\textwidth}
\footnotesize

$^{a}$ Talos on-chip energy excludes the external DRAM contribution,
i.e.,
$E_{\mathrm{on-chip}} =
E_{\mathrm{workload}} - E_{\mathrm{DRAM}}$.
The corresponding on-chip power is computed over the complete
workload execution.

$^{b}$ Total Talos energy and power include the modeled external DRAM
contribution. The total workload energy is also the quantity used by
the optimization when energy is included as an objective.

$^{c}$ SAURIA does not report a directly comparable VGG-16
energy-per-inference value including the external DRAM contribution.
Consequently, neither total energy nor total system power is shown.

$^{d}$ SAURIA reports a post-layout FP16 implementation in
GF22FDX 22\,nm with an area of 0.897\,mm$^2$, an average accelerator
power of 258\,mW at peak throughput, and a peak throughput of
284\,GFLOP/s \cite{sauria}. The reported power corresponds to the
accelerator and is therefore compared with the Talos on-chip power
rather than its total power including DRAM.

The comparison should be interpreted as indicative rather than as a
direct benchmark. Talos and SAURIA use different technology nodes,
and the published SAURIA performance is reported as peak arithmetic
throughput rather than VGG-16 inferences per second.

\end{minipage}
\end{table*}

\begin{table*}[t]
\centering
\caption{Comparison of the Talos INT16 AlexNet convolutional
configurations selected under different optimization objectives and
the Eyeriss v1 accelerator.}
\label{tab:eyeriss_v1_comparison}

\setlength{\tabcolsep}{3.2pt}
\renewcommand{\arraystretch}{1.15}

\resizebox{0.99\textwidth}{!}{%
\begin{tabular}{@{}lcccccccc@{}}
\hline
\textbf{Configuration} &
\textbf{Technology} &
\textbf{Precision} &
\makecell{\textbf{On-chip Energy}\\\textbf{(mJ/inf.)$^{a}$}} &
\makecell{\textbf{Energy incl. DRAM}\\\textbf{(mJ/inf.)$^{b}$}} &
\textbf{Area} &
\makecell{\textbf{On-chip}\\\textbf{Power$^{a}$}} &
\makecell{\textbf{Total}\\\textbf{Power$^{b}$}} &
\makecell{\textbf{Latency}\\\textbf{(cycles)$^{c}$}} \\
\hline

Talos (E) &
65\,nm CMOS &
INT16 &
44.060 &
57.059 &
3.251\,mm$^2$ &
0.728\,W &
0.943\,W &
12107646 \\

Talos (E+A) &
65\,nm CMOS &
INT16 &
56.486 &
71.998 &
2.828\,mm$^2$ &
0.616\,W &
0.785\,W &
18336893 \\

Talos (A) &
65\,nm CMOS &
INT16 &
45.536 &
214.718 &
0.265\,mm$^2$ &
0.058\,W &
0.275\,W &
155908379 \\

Talos (P) &
65\,nm CMOS &
INT16 &
550.904 &
553.165 &
386.255\,mm$^2$ &
74.224\,W &
74.528\,W &
1484444 \\

Talos (A+P) &
65\,nm CMOS &
INT16 &
44.060 &
57.059 &
3.251\,mm$^2$ &
0.728\,W &
0.943\,W &
12107646 \\

Talos (E+P) &
65\,nm CMOS &
INT16 &
78.835 &
87.571 &
46.633\,mm$^2$ &
10.442\,W &
11.599\,W &
1509958 \\

Talos (E+A+P) &
65\,nm CMOS &
INT16 &
44.060 &
57.059 &
3.251\,mm$^2$ &
0.728\,W &
0.943\,W &
12107646 \\

\hline

Eyeriss v1 \cite{chen2017eyeriss} &
65\,nm CMOS &
INT16 &
8.01$^{d}$ &
--$^{e}$ &
12.25\,mm$^2$ &
0.278\,W$^{d}$ &
--$^{e}$ &
5763689$^{f}$ \\

\hline
\end{tabular}%
}

\vspace{2mm}
\begin{minipage}{0.99\textwidth}
\footnotesize

$^{a}$ Talos on-chip figures include the selected PEs and on-chip
memory hierarchy, excluding the external DRAM contribution.

$^{b}$ Talos total figures include the external DRAM power and energy
estimated from the workload activity.

$^{c}$ The Talos workload reproduces the Eyeriss AlexNet convolutional
benchmark with batch size 4. Energy and latency are normalized per
image. Both Talos and Eyeriss v1 operate at a 200\,MHz core frequency
in this comparison.

$^{d}$ Eyeriss v1 reports a measured chip power of 278\,mW and a
throughput of 34.7 frames/s for the five convolutional layers of
AlexNet at 1\,V. Its energy per image is therefore derived as
$E_{\mathrm{inf}}=P/(\mathrm{inf/s})\approx8.01$\,mJ/inf.

$^{e}$ Eyeriss v1 reports the external DRAM traffic (15.4\,MB for a
batch of four AlexNet images), but does not report the corresponding
external-memory energy or total system power. These quantities are
therefore left unspecified.

$^{f}$ Eyeriss v1 latency is expressed as equivalent cycles per image
and is derived from its measured 34.7 frames/s throughput at 200\,MHz:
$N_{\mathrm{cycles}}=200\times10^6/34.7\approx5.76\times10^6$ cycles.

\end{minipage}

\end{table*}

\begin{table*}[t]
\centering
\caption{Comparison of the Talos INT8 MobileNet V1 0.5/128
configurations selected under different optimization objectives and
the Eyeriss v2 accelerator.}
\label{tab:eyerissv2_mobilenet_comparison}

\setlength{\tabcolsep}{3.2pt}
\renewcommand{\arraystretch}{1.15}

\resizebox{0.99\textwidth}{!}{%
\begin{tabular}{@{}lcccccccc@{}}
\hline
\textbf{Configuration} &
\textbf{Technology} &
\textbf{Precision} &
\makecell{\textbf{On-chip Energy}\\\textbf{(mJ/inf.)$^{a}$}} &
\makecell{\textbf{Energy incl. DRAM}\\\textbf{(mJ/inf.)$^{b}$}} &
\textbf{Area} &
\makecell{\textbf{On-chip}\\\textbf{Power$^{a}$}} &
\makecell{\textbf{Total}\\\textbf{Power$^{b}$}} &
\makecell{\textbf{Latency}\\\textbf{(cycles)$^{f}$}} \\
\hline

Talos (E) &
65\,nm CMOS &
INT8 &
0.972 &
2.795 &
5.672\,mm$^2$ &
1.199\,W &
3.450\,W &
162026 \\

Talos (E+A) &
65\,nm CMOS &
INT8 &
2.881 &
7.276 &
0.441\,mm$^2$ &
0.090\,W &
0.227\,W &
6416591 \\

Talos (A) &
65\,nm CMOS &
INT8 &
1.889 &
17.320 &
0.170\,mm$^2$ &
0.041\,W &
0.373\,W &
9284696 \\

Talos (P) &
65\,nm CMOS &
INT8 &
14.094 &
15.524 &
110.515\,mm$^2$ &
24.326\,W &
26.793\,W &
115879 \\

Talos (A+P) &
65\,nm CMOS &
INT8 &
1.399 &
4.828 &
1.536\,mm$^2$ &
0.388\,W &
1.338\,W &
721433 \\

Talos (E+P) &
65\,nm CMOS &
INT8 &
2.079 &
3.540 &
15.237\,mm$^2$ &
3.097\,W &
5.275\,W &
134229 \\

Talos (E+A+P) &
65\,nm CMOS &
INT8 &
1.388 &
4.816 &
1.536\,mm$^2$ &
0.385\,W &
1.335\,W &
721433 \\

\hline

Eyeriss v2 dense \cite{chen2019eyerissv2} &
65\,nm TSMC LP &
INT8 &
0.508$^{c}$ &
--$^{d}$ &
--$^{e}$ &
0.651\,W$^{c}$ &
--$^{d}$ &
155994$^{g}$ \\

\hline
\end{tabular}%
}

\vspace{2mm}
\begin{minipage}{0.99\textwidth}
\footnotesize

$^{a}$ Talos on-chip figures include the selected PEs and on-chip
memory hierarchy, excluding the external DRAM contribution.

$^{b}$ Talos total figures include the modeled external DRAM
contribution. The total energy is also the quantity used by the
optimization whenever energy is included as an objective.

$^{c}$ Eyeriss v2 reports inference/J and inference/s rather than
energy per inference and average power directly. For the dense
MobileNet V1 0.5/128 workload, the paper reports 1969.8 inf/J and
1282.1 inf/s. The values shown here are therefore derived as
\[
E_{\mathrm{inf}} =
\frac{1}{1969.8}
\approx 0.508\,\text{mJ/inf.},
\]
and
\[
P =
E_{\mathrm{inf}}\times 1282.1
\approx 0.651\,\text{W}.
\]
These values are treated as on-chip figures because the paper reports
external DRAM traffic separately and its power breakdown is given for
the accelerator's on-chip components.

$^{d}$ The Eyeriss v2 paper reports external DRAM traffic for the
MobileNet workload, but does not provide the corresponding DRAM energy
or total system power. Therefore, no directly comparable value is
reported for energy including DRAM or total power.

$^{e}$ Eyeriss v2 is reported as an implementation with 2695k
NAND-2-equivalent logic gates and 246\,kB of on-chip SRAM rather than
as a directly comparable silicon area in mm$^2$. The area is therefore
left unspecified in this table.

$^{f}$ Talos and Eyeriss v2 operate at 200\,MHz in this comparison.
Talos latency is reported directly in workload cycles at this
frequency.

$^{g}$ Eyeriss v2 latency in cycles is derived from the reported
throughput at 200\,MHz:
\[
N_{\mathrm{cycles}} =
\frac{200\times10^6}{1282.1}
\approx 155994
\]
cycles per inference. The Eyeriss v2 benchmark uses batch size 1,
MobileNet V1 with width multiplier 0.5, and a $128\times128$ input.

The dense Eyeriss v2 result is used as the direct reference because
the Talos workload models a dense MobileNet V1 0.5/128 network and
does not reproduce the sparsity exploited by the sparse Eyeriss v2
configuration.

\end{minipage}

\end{table*}
%\subsection{Impact of the Hierarchical Exploration}


%\subsubsection{Level 1 and Level 2 Ranking Comparison}

% One of the main thesis results.
%
% Compare architectural preference before and after physical IP
% instantiation.
%
% Important questions:
% Does the best L1 point remain the best after IP selection?
% Can an apparently worse L1 architecture obtain a better physical
% implementation?
%
% Do this from the new final runs.


%\subsubsection{Effect of Hardware Constraints}

% Final constraint experiments.
%
% Suggested cases:
% no constraints
% area constraint
% power constraint
% minimum-frequency constraint
% combined constraints.
%
% Repeat with one characterized pool and one representative workload.


%\subsubsection{Design-Space Reduction}

% Compare:
% theoretical L1 architecture space
% architectures actually evaluated
% architectures transferred to L2
% theoretical/compatible IP combinations
% L2 combinations actually evaluated.
%
% This is the direct experiment for the hierarchical-search
% hypothesis.


%\subsubsection{Exploration Runtime}

% Measure using the new final runs.
%
% L1 time
% L2 time
% total time
% characterization time if relevant
%
% If possible, run the same small L2 space using exhaustive and
% NSGA-II.


%\subsection{Workload Case Studies}

%\subsubsection{Integer Workload}

% Use the characterized integer pool.
%
% Candidate workloads:
% SqueezeNet INT8
% AlexNet INT16, with the precision compatibility explicitly
% considered.
%
% Report final architectural and physical alternatives.


%\subsubsection{Floating-Point Workload}

% Use FP16 and/or FP32 pool.
%
% Candidate workloads:
% SqueezeNet FP16
% ResNet18 first layer FP32.
%
% Keep workload and PE precision coherent for the final quantitative
% comparison.


%\subsubsection{Cross-Workload Comparison}

% Compare the final results from at least two networks.
%
% Questions:
% Does the preferred L1 architecture change?
% Does the preferred PE change?
% Does memory selection change?
% Are the physical constraints equally important for each workload?
%
% Do not compare incompatible workload/PE precisions as if they were
% equivalent implementations.


\subsection{Results Summary}

% Fill only after all final experiments have been rerun.
%
% Keep quantitative:
% design-space reduction
% ranking changes
% constraint effect
% selected IP changes
% workload dependence.
